\exercisegroup

\begin{enumerate}
\def\labelenumi{\arabic{enumi}.}
\tightlist
\item
  \begin{enumerate}
  \def\labelenumii{(\alph{enumii})}
  \tightlist
  \item
    设 E 是一个集合，并设 \(\mathfrak{F}(\mathbf{E})\) 是 E 的有限子集的集合。证明 \(\mathfrak{F}(\mathbf{E})\) 是满足以下条件的最小子集 \(\mathfrak{G}\) （\(\mathfrak{P}(\mathbf{E})\) 中的子集）：(i) \(\phi \in \mathfrak{G}\) ；(ii) 关系 \(\mathbf{X} \in \mathfrak{G}\) 和 \(x \in \mathbf{E}\) 蕴含
  \end{enumerate}
\end{enumerate}

\[
\mathbf {X} \cup \{x \} \in \mathfrak {G}.
\]

\begin{enumerate}
\def\labelenumi{(\alph{enumi})}
\setcounter{enumi}{1}
\item
  由(a)推出，E的两个有限子集A、B的并集是有限的（考虑E的子集X的集合，使得 \(X \cup A\) 是有限的；参见§5，第1号，命题1，推论1）。
\item
  从(a)和(b)推出，对于每个有限集E，集合 \(\mathfrak{P}(E)\) 是有限的（考虑E的子集X的集合，使得 \(\mathfrak{P}(X)\) 是有限的；参见§5第1号命题1推论4）。
\end{enumerate}

\begin{enumerate}
\def\labelenumi{\arabic{enumi}.}
\setcounter{enumi}{1}
\item
  证明集合 E 是有限的当且仅当 \(\mathfrak{P}(\mathbf{E})\) 的每个非空子集（关于包含关系）都有极大元。（为证明该条件是充分的，将其应用于集合 \(\mathfrak{F}(\mathbf{E})\) （即E的有限子集的集合。）
\item
  证明：若良序集E在赋予相反序关系后所得的序集也是良序的，则E是有限的（考虑E的最大元x，使得段 \(S_{x}\) 是有限的）。
\item
  设 E 是一个有限集，具有 \(n \geqslant 2\) 个元素，并令 C 为 \(E \times E\) 的一个子集，使得对于E中每一对不同的元素x、y，两个元素 \((x, y)\) , \((y, x)\) （\(E \times E\) 的元素）中恰好有一个属于 C。证明存在一个区间 {[}1, n{]} 到 E 上的映射 f，使得 \((f(i), f(i+1)) \in \mathbb{C}\) 对于 \(1 \leqslant i \leqslant n - 1\) （对 n 使用归纳法）。
\end{enumerate}

¶ 5. 设 E 是一个有序集，且存在一个整数 k，使得 k 是 E 的自由子集 X（§ 1，习题 5）中元素个数的最大值。证明 E 可以划分为 k 个全序子集（关于诱导序）。证明分两步：

\begin{enumerate}
\def\labelenumi{(\alph{enumi})}
\item
  若 E 是有限的且有 n 个元素，则对 n 使用归纳法；设 a 是 E 的一个极小元素，并令 \(E' = E - \{a\}\) 。如果存在一个将 \(E'\) 划分为 k 个全序集 \(C_i (1 \leqslant i \leqslant k)\) ，令 \(U_i\) 为所有 \(x \in C_i\) （它们 \(\geqslant a\) ）所成的集合 . 证明存在至少一个指标 i，使得 \(E' - U_i\) 的一个自由子集至多有 k - 1 个元素。其证明采用归谬法。对每个 i，设 \(S_i\) 是 \(E' - U_i\) 的一个自由子集，且具有 k 个元素，设 S 为这些集合 \(S_i\) 的并集，并且令 \(s_j\) 为 \(S \cap C_j\) 的最小元素（对每个指标 \(j \leqslant k\) ）；证明 \(k + 1\) 个元素 \(a, s_1, \ldots, s_k\) 构成E的一个自由子集。
\item
  如果E是任意的，则通过对k进行归纳来证明，如下。E的子集C被称为在E中强相关的，如果对于E的每个有限子集F，存在F的一个划分，使得该划分由至多k个全序集组成，并且 \(C \cap F\) 包含于其中之一。证明存在一个极大的强相关子集 \(C_{0}\) ，并且 \(E - C_{0}\) 的每个自由子集至多有 k - 1 个元素。（用反证法论证，假设存在一个自由子集 \(\{a_{1}, \ldots, a_{k}\}\) （具有 k 个元素，在 \(E - C_{0}\) 中）。考虑每个集合 \(C_{0} \cup \{a_{i}\} (1 \leqslant i \leqslant k)\) ，并表达它并非强相关这一事实，从而为每个指标 i 引入 E 的一个有限子集 \(F_{i}\) 。然后考虑这些集合 \(F_{i}\) 的并集 F，并利用 \(C_{0}\) 是强相关的事实来得出矛盾。）
\end{enumerate}

\begin{enumerate}
\def\labelenumi{\arabic{enumi}.}
\setcounter{enumi}{5}
\tightlist
\item
  \begin{enumerate}
  \def\labelenumii{(\alph{enumii})}
  \tightlist
  \item
    设 A 为一个集合，且设 \((\mathbf{X}_{i})_{1\leqslant i\leqslant m},(\mathbf{Y}_{j})_{m+1\leqslant j\leqslant m+n}\) 是 A 的子集的两个有限族。设 h 为满足以下条件的最小整数：对每个整数 \(r\leqslant m-h\) 以及每个子集 \(\{i_{1},\ldots,i_{r+h}\}\) （\(r+h\) 个元素，取自 {[}1, m{]}），存在一个子集 \(\{j_{1},\ldots,j_{r}\}\) ，它由 \([m+1,m+n]\) 中的r个元素组成，使得诸集合 \(X_{i_{\alpha}}(1\leqslant\alpha\leqslant r+h)\) 的并集与每一个集合 \(Y_{j_{3}}(1\leqslant\beta\leqslant r)\) 相交（这意味着 \(m\leqslant n+h\) ). 证明存在 A 的一个有限子集 B，其元素个数至多为 \(n+h\) 个，使得每个 \(X_{i}(1\leqslant i\leqslant m)\) 以及每一个 \(Y_{j}(m+1\leqslant j\leqslant m+n)\) 都与 B 相交。（考虑区间 \([1,m+n]\) 上的序关系，其图形是对角线与由点对 \((i,j)\) （使得 \(1\leqslant i\leqslant m\) 、 \(m+1\leqslant j\leqslant m+n\) 且 \(X_{i}\cap Y_{j}\neq\varnothing\) ）所组成的集合的并集，并将习题5应用于这个有序集。）
  \end{enumerate}
\end{enumerate}

\begin{enumerate}
\def\labelenumi{(\alph{enumi})}
\setcounter{enumi}{1}
\item
  让 \(\mathbf{E}\) 和 \(\mathbf{F}\) 是两个有限集，且设 \(x \to \mathbf{A}(x)\) 为从 \(\mathbf{E}\) 到 \(\mathfrak{P}(\mathbf{F})\) 中的一个映射。则存在一个单射 \(f\) （从 \(\mathbf{E}\) 到 \(\mathbf{F}\) 中），使得 \(f(x) \in \mathbf{A}(x)\) （对每一个 \(x \in \mathbf{E}\) ），当且仅当对 \(\mathbf{E}\) 的每个子集 \(\mathbf{H}\) 我们拥有 \(\operatorname{Card}\left(\bigcup_{x \in \mathbf{H}} \mathbf{A}(x)\right) \geqslant \operatorname{Card}(\mathbf{H})\) （证明方法类似于（a）的证明，其中 \(h = 0\) ).
\item
  在(b)的假设下，设 \(\mathbf{G}\) 为 \(\mathbf{F}\) 的一个子集。则存在一个单射 \(f\) （从 \(\mathbf{E}\) 到 \(\mathbf{F}\) 中），使得 \(f(x) \in \mathbf{A}(x)\) （对每一个 \(x \in \mathbf{E}\) ），并且使得 \(f(\mathbf{E}) \supset \mathbf{G}\) ，当且仅当 \(f\) 满足(b)的条件，并且对每个子集 \(\mathbf{L}\) （\(\mathbf{G}\) 的子集），所有 \(x \in \mathbf{E}\) （满足 \(\mathbf{A}(x) \cap \mathbf{L} \neq \emptyset\) ）所组成的集合的基数是 \(\geqslant \operatorname{Card}(\mathbf{L})\) 。（令 \((a_i)_{1 \leqslant i \leqslant p}\) 为 \(\mathbf{G}\) 的不同元素的序列，按某种顺序排列；令 \((b_j)_{p+1 \leqslant j \leqslant p+m}\) 为 \(\mathbf{F}\) 的不同元素的序列，按某种顺序排列；并令
\end{enumerate}

\[
\left(c _ {k}\right) _ {p + m + 1 \leqslant k \leqslant p + m + n}
\]

为E的不同元素的序列，按某种顺序排列。考虑集合 \([1, p + m + n]\) 上的序关系，其图形是对角线与由满足以下条件的点对 \((i, j)\) 所组成的集合的并集：要么

\[
1 \leqslant i \leqslant p \quad \text { 且 } \quad p + 1 \leqslant j \leqslant p + m \quad \text { 且 } \quad a _ {i} = b _ {j},
\]

\[
\text { 或 } \quad 1 \leqslant i \leqslant p \quad \text { 且 } \quad p + m + 1 \leqslant j \leqslant p + m + n \quad \text { 且 } \quad a _ {i} \in A (c _ {j}),
\]

要么 \(p + 1 \leqslant i \leqslant p + m\) 且 \(p + m + 1 \leqslant j \leqslant p + m + n\) 且 \(b_i \in \mathrm{A}(c_j)\) ;

然后应用练习5。）

\begin{enumerate}
\def\labelenumi{\arabic{enumi}.}
\setcounter{enumi}{6}
\tightlist
\item
  一个元素 \(a\) （格 \(\mathbf{E}\) 的元素）被称为不可约的，如果关系 \(\sup (x, y) = a\) 蕴含要么 \(x = a\) 要么 \(y = a\) .
\end{enumerate}

\begin{enumerate}
\def\labelenumi{(\alph{enumi})}
\item
  证明在有限格 \(\mathbf{E}\) 中，每个元素 \(a\) 可以写成 \(\sup (e_1, \ldots, e_k)\) ，其中 \(e_i (1 \leqslant i \leqslant k)\) 都是不可约的。
\item
  设E是一个有限格，J为其不可约元素的集合。对每个 \(x \in E\) 令 \(S(x)\) 为所有 \(y \in J\) （它们 \(\leqslant x\) ）所成的集合。证明映射 \(x \to S(x)\) 是E到 \(\mathfrak{P}(J)\) （按包含关系排序）的一个子集上的同构，并且 \(S(\inf(x, y)) = S(x) \cap S(y)\) .
\end{enumerate}

\begin{enumerate}
\def\labelenumi{\arabic{enumi}.}
\setcounter{enumi}{7}
\tightlist
\item
  \begin{enumerate}
  \def\labelenumii{(\alph{enumii})}
  \tightlist
  \item
    设 E 是一个分配格（§ 1，习题 16）。如果 \(a\) 在E中不可约（习题7），证明该关系 \(a \leqslant \sup (x, y)\) 蕴含 \(a \leqslant x\) 或 \(a \leqslant y\) .
  \end{enumerate}
\end{enumerate}

\begin{enumerate}
\def\labelenumi{(\alph{enumi})}
\setcounter{enumi}{1}
\item
  设 E 是一个有限分配格，并设 J 为其不可约元素的集合，按诱导序排列。证明该同构 \(x \to \mathrm{S}(x)\) （将 E 映到 \(\mathfrak{P}(\mathbf{J})\) 的一个子集上，定义于练习7（b）中）满足 \(\mathrm{S}(\sup(x,y)) = \mathrm{S}(x) \cup \mathrm{S}(y)\) 。由此推出，若 \(J^{*}\) 是将 J 赋予相反序关系后得到的有序集，则 E 同构于递增映射所组成的集合 \(\mathcal{A}(\mathbf{J}^{*}, \mathbf{I})\) （由 \(J^{*}\) 到 \(I = \{0, 1\}\) 中的递增映射）（§1，习题6）。
\item
  在（b）的假设下，设 P 为 J 中除 E 的最小元素以外的元素组成的集合。对每个 \(x \in E\) ，让 \(y_{1}, \ldots, y_{k}\) 为区间 \(]x, \rightarrow[\) 在 E 中的互不相同的极小元素；对每个指标 i，设 \(q_{i}\) 为 P 中的一个元素，使得 \(q_{i} \notin S(x)\) 和 \(q_{i} \in S(y_{i})\) 。证明元素 \(q_{1}, \ldots, q_{k}\) 中任意两个都不可比较。
\end{enumerate}

\begin{enumerate}
\def\labelenumi{(\alph{enumi})}
\setcounter{enumi}{3}
\tightlist
\item
  反之，设 \(q_1, \ldots, q_k\) 是 \(k\) 个元素（属于 \(\mathbf{P}\) ），其中任意两个都不可比较。设 \(u = \sup (q_1, \ldots, q_k)\) 并令
\end{enumerate}

\[
v _ {i} = \sup _ {1 \leqslant j \leqslant k, j \neq i} (q _ {j}) \quad (1 \leqslant i \leqslant k).
\]

证明 \(v_{i} < u\) 对于 \(1 \leqslant i \leqslant k\) . 设 \(x = \inf (v_{1}, \ldots, v_{k})\) 并令

\[
y _ {i} = \inf _ {1 \leqslant j \leqslant k, J \neq i} (v _ {j}).
\]

证明 \(x < y_i\) 对于每个指标 \(i\) ，并推断出区间 \(]x, \to[\) 至少有 \(k\) 个不同的极小元素。

\begin{enumerate}
\def\labelenumi{\arabic{enumi}.}
\setcounter{enumi}{8}
\tightlist
\item
  格E的子集A被称为子格，如果对于 A 的元素的每一对 \((x, y)\) ， \(\sup_{\mathbf{E}}(x, y)\) 和 \(\inf_{\mathbf{E}}(x, y)\) 属于A。
\end{enumerate}

\begin{enumerate}
\def\labelenumi{(\alph{enumi})}
\tightlist
\item
  设 \((\mathbf{C}_i)_{1\leqslant i\leqslant n}\) 为一个全序集的有限族，并令
\end{enumerate}

\[
\mathrm{E} = \prod_ {i = 1} ^ {n} \mathrm{C} _ {i}
\]

设它们的乘积为 E。令 A 为 E 的一个子格。证明 A 不能有多于 n 个两两不可比较的不可约元素（习题 7）。（证明采用归谬法。假设存在 r \textgreater{} n 个不可约元素 \(a_{1}, \ldots, a_{r}\) 在 A 中，且其中任意两个都不可比较。考虑 \(u = \sup(a_{1}, \ldots, a_{r})\) 并且

\[
v _ {j} = \sup _ {\mathbf {1} \leqslant j \leqslant r, j \neq i} (a _ {j})
\]

的元素。通过向各因子投影，证明 \(u = v_{i}\) 对某个指标 i 成立，从而诸 \(a_{i}\) 中有两个是可比较的。）

\begin{enumerate}
\def\labelenumi{(\alph{enumi})}
\setcounter{enumi}{1}
\tightlist
\item
  反之，设 F 是一个有限分配格，设 P 是 F 中除 F 的最小元素以外的不可约元素的集合，并设 n 是 P 的自由子集中元素个数的最大值（§1，习题 5）。证明 F 同构于 n 个全序集的乘积的一个子格。（应用习题 5，它表明 P 是 n 个两两无公共元素的全序集 \(P_{i}\) 的并。设 \(C_{i}\) 为通过向 \(P_{i}\) 添加一个最小元素而得到的全序集 ( \(1 \leqslant i \leqslant n\) ）。对于每个 \(x \in F\) ，关联族 \((x_{i})_{1 \leqslant i \leqslant n}\) ，其中 \(x_{i}\) 是 \(C_{i}\) 中由 \(P_{i}\) 的元素（它们 \(\leqslant x\) ）所组成的集合的上确界 .
\end{enumerate}

¶ 10. (a) 一个有序集 E 同构于 n 个全序集的乘积的一个子集，当且仅当 E 上的序关系的图是 E 上 n 个全序关系的图的交集。（为证明该条件是必要的，证明若

\[
\mathbf {F} = \prod_ {i = 1} ^ {n} \mathbf {F} _ {i}
\]

是 n 个全序集的乘积，则 F 上乘积序的图是 F 上 n 个字典序的图的交集。）

\begin{enumerate}
\def\labelenumi{(\alph{enumi})}
\setcounter{enumi}{1}
\item
  一个有序集 E 同构于两个全序集的乘积的某个子集，当且仅当该序 \(\Gamma\) （E 上的）使得存在另一个序 \(\Gamma'\) （E 上的），且具有如下性质：E 中任意两个不同元素恰好关于这两个序 \(\Gamma\) , \(\Gamma'\) 中的一个序是可比较的 .
\item
  设 A 为含 n 个元素的有限集。设 E 为 \(\mathfrak{P}(A)\) 中由所有子集 \(\{x\}\) 以及 A--- \(\{x\}\) （当x遍历A时）组成的子集。证明n是使得按包含关系排序的E同构于m个全序集的乘积的子集的最小整数m（利用(a)）。
\end{enumerate}

\begin{enumerate}
\def\labelenumi{\arabic{enumi}.}
\setcounter{enumi}{10}
\tightlist
\item
  设A为一个集合，并设R为A的有限子集集合F(A)的一个子集。若R满足以下条件，则称R是可移动的：
\end{enumerate}

（MO）如果 X、Y 是 R 的两个不同元素，并且如果 \(z \in X \cap Y\) ，那么存在 \(Z \subset X \cap Y\) 属于 R 且满足 \(z \notin Z\) .

A 的子集 P 被称为纯的，如果它不包含属于 R 的任何集合。

\begin{enumerate}
\def\labelenumi{(\alph{enumi})}
\item
  证明 \(\mathbf{A}\) 的每一个纯子集都包含在 \(\mathbf{A}\) 的一个极大纯子集中 .
\item
  设 M 是 A 的一个极大纯子集。证明对每个 \(x \in \complement M\) 存在 M 的唯一的有限子集 \(\mathbf{E}_{\mathbf{M}}(x)\) ，使得 \(\mathbf{E}_{\mathbf{M}}(x) \cup \{x\} \in \Re\) 。此外，如果 \(y \in \mathbf{E}_{\mathbf{M}}(x)\) ，集合 \((\mathbf{M} \cup \{x\}) - \{y\}\) 是 A 的一个极大纯子集。
\item
  设M、N是A的两个极大纯子集，使得 \(N \cap GM\) 是有限的。证明 Card (M) = Card (N)。（对 \(N \cap GM\) 的基数进行归纳证明，利用（b）。）
\item
  设 M, N 为 A 的两个极大纯子集，并令 \(N' = N \cap CM\) , \(M' = M \cap CN\) . 证明
\end{enumerate}

\[
\mathbf {M} ^ {\prime} \subset \bigcap_ {x \in \mathbb {N} ^ {\prime}} \mathrm{E} _ {\mathbf {M}} (x).
\]

* 由此推出 \(\operatorname{Card}(\mathbf{M}) = \operatorname{Card}(\mathbf{N})\) （根据 (c)，我们可归结为 \(\mathbf{N}'\) 和 \(\mathbf{M}'\) 为无限的情形；然后证明 \(\operatorname{Card}(\mathbf{M}') \leqslant \operatorname{Card}(\mathbf{N}')\) ). *

